\chapter{Introduction}
\label{ch:intro}

\section{Motivation}
In a full classroom the \cotwo{} that the people breathe out builds up quickly. High \cotwo{} levels are
a sign of stale air and are linked to tiredness and poorer concentration \todo{reference on CO2 and
cognitive performance}. The usual answer is to open the windows. This works, but in winter every open
window lets heat out, and since most rooms have no sensors, nobody knows how long and how far the windows
need to be open.

A model that forecasts the \cotwo{} in the room a few minutes ahead could help with this decision. With
such a forecast, the windows could be opened only when the \cotwo{} is about to become too high, and
closed again as soon as it is safe. Computing the airflow in the room with computational fluid dynamics
(CFD) is accurate, but one simulation takes hours, which is far too slow for a decision every few minutes.
\citet{bian2025} showed that a neural operator trained on CFD results can forecast the \cotwo{} in a
classroom within a fraction of a second and can be used for control.

\section{Goal of this thesis}
The goal of this thesis is to apply the method of \citet{bian2025} to a real classroom of our university,
which is ventilated through windows instead of ceiling vents, and to prepare it for the real sensor data of
the wireless sensor network (WSN) that is being installed in the room. In particular I want to answer
these questions:
\begin{enumerate}
  \item Can the method be transferred to our room with open-source software (OpenFOAM instead of
        ANSYS Fluent)?
  \item How accurate is the generated training data, and how accurate is the forecast?
  \item Can the simulations be made faster, for example by running them on the GPU?
  \item Can the network be trained from the equations alone, without simulation data?
  \item What is needed to use the model with the real sensor data?
\end{enumerate}

\section{Approach}
\Cref{fig:pipeline} shows the pipeline. For every window setting I simulate the airflow in the room with
OpenFOAM. From this airflow my own GPU solver computes the \cotwo{} of the people in the room for
30~minutes. The forecast model, a GNOT-type neural operator, is trained on these results: it receives the
\cotwo{} at 1.6~m height over the last 6~minutes, the window speeds and the number of people, and predicts
the \cotwo{} for the next 3~minutes.

\begin{figure}[h]
  \centering
  \figplaceholder{pipeline: room geometry $\rightarrow$ OpenFOAM airflow $\rightarrow$ own \cotwo{} solver
  (GPU) $\rightarrow$ dataset $\rightarrow$ GNOT training $\rightarrow$ forecast}
  \caption{Overview of the pipeline. \todo{draw}}
  \label{fig:pipeline}
\end{figure}

\section{Contributions}
\begin{itemize}
  \item A complete, scripted pipeline from the room geometry to the trained forecast model, built only on
        open-source software.
  \item A mass-conserving GPU solver for the \cotwo{} that reproduces OpenFOAM's own \cotwo{} within about
        2\,\%, and a study of how often the airflow has to be saved to reach this accuracy.
  \item A corrected dataset for the current room setup and a comparison of models trained on the old and the
        corrected data.
  \item A test of GPU alternatives to OpenFOAM, including an own lattice Boltzmann solver, and a survey of
        the available software.
  \item A test of training without simulation data (physics-informed), which shows the limits of this idea
        for our room.
\end{itemize}

\section{Structure}
\Cref{ch:background} gives the background on CFD, turbulence modelling and neural operators.
\Cref{ch:room} describes the room and the OpenFOAM simulations, and \cref{ch:dataset} how the \cotwo{}
dataset is built and checked. \Cref{ch:model} describes the forecast model and \cref{ch:results} its
results. \Cref{ch:computation} covers run times and the GPU alternatives, \cref{ch:pinn} the training
without data, and \cref{ch:sensors} the preparation for the real sensor data. \Cref{ch:conclusion}
summarises the results.
